\section{Concluding Remarks and Future Direction}
\label{sec:conclusion}
Our work was motivated by the need to understand the key considerations and design goals that impact human experiences within smart buildings. To achieve this, we conducted a scoping review of 192 papers, adopting a hybrid inductive-deductive approach that draws on thematic analysis coding techniques \cite{braun2023thematic}. Through this process, we identified 11 distinct human experiences, 20 design mechanisms that shape these experiences as well as developed a typology of the technological interventions through which they are implemented.

Our research sheds light on how current smart building designs aim to address human experiences. However, this is only a partial view of the broader challenge \cite{kirsh2019architects}. While we have explored \emph{what} we design for on a human level, an even more pressing question remains: what \emph{should} we design for? To fully answer this question, we need to engage deeply with both the users who inhabit smart buildings and the practitioners who design them. Looking ahead, answering the question of what \emph{should} we design for entails more than just academic inquiry. It demands a broader conversation involving interdisciplinary collaboration, policy-making, and the inclusion of diverse voices. This scoping review offers a foundation for understanding the current state of smart building design and research, but it also opens the door to more pressing inquiries. Moving forward, the challenge will be to create practical opportunities for smart building development underpinned by a design horizon that prioritises the voices and visions of those who inhabit them. 